% =====================================================================
%  capitulos/00b_nomenclatura.tex  --  Lista de simbolos y abreviaturas
%  Material preliminar. Se incluye desde main.tex despues de \listoftables
%  y antes de \pagenumbering{arabic}. Tabla manual (longtable) para no
%  depender del paquete nomencl.
% =====================================================================

\chapter*{Nomenclatura}
\addcontentsline{toc}{chapter}{Nomenclatura}
\markboth{Nomenclatura}{Nomenclatura}

Se reúnen los símbolos y las abreviaturas empleados en el documento. Las
unidades se indican en el Sistema Internacional como referencia dimensional;
en los casos de validación se emplean además sistemas técnicos coherentes
(por ejemplo, kgf/cm), según se aclara en cada problema.

\section*{Símbolos}

\renewcommand{\arraystretch}{1.2}
\begin{longtable}{@{}p{2.4cm}p{8.4cm}p{3.0cm}@{}}
\toprule
\textbf{Símbolo} & \textbf{Significado} & \textbf{Unidad} \\
\midrule
\endfirsthead
\toprule
\textbf{Símbolo} & \textbf{Significado} & \textbf{Unidad} \\
\midrule
\endhead
\bottomrule
\endfoot
$\Omega,\ \Gamma$ & Dominio del problema y su frontera & --- \\
$\Gamma_u,\ \Gamma_t$ & Frontera con desplazamientos / tracciones prescritos & --- \\
$(x,y)$ & Coordenadas físicas (cartesianas) & m \\
$(\xi,\eta)$ & Coordenadas naturales del elemento de referencia & adimensional \\
$E$ & Módulo de elasticidad (de Young) & Pa \\
$\nu$ & Coeficiente de Poisson & adimensional \\
$\rho$ & Densidad del material & kg/m$^3$ \\
$t$ & Espesor del elemento & m \\
$h$ & Tamaño característico de malla & m \\
$p$ & Orden polinómico del elemento & adimensional \\
$\bm{u}$ & Vector de desplazamientos nodales & m \\
$u,\ v$ & Componentes de desplazamiento en $x$ e $y$ & m \\
$\bm{b}$ & Fuerza por unidad de volumen & N/m$^3$ \\
$\bm{t}$ & Vector de tracciones de borde & N/m$^2$ \\
$\bm{\sigma}$ & Vector de tensiones $[\sigma_x,\sigma_y,\tau_{xy}]^T$ (Voigt) & Pa \\
$\sigma_x,\ \sigma_y$ & Tensiones normales & Pa \\
$\tau_{xy}$ & Tensión cortante & Pa \\
$\sigma_1,\ \sigma_2$ & Tensiones principales & Pa \\
$\sigma_{\mathrm{VM}}$ & Tensión equivalente de von Mises & Pa \\
$\bm{\varepsilon}$ & Vector de deformaciones $[\varepsilon_x,\varepsilon_y,\gamma_{xy}]^T$ & adimensional \\
$\varepsilon_x,\ \varepsilon_y$ & Deformaciones normales & adimensional \\
$\gamma_{xy}$ & Deformación angular (de corte) & adimensional \\
$N_i$ & Función de forma asociada al nodo $i$ & adimensional \\
$\bm{N}$ & Matriz de funciones de forma & adimensional \\
$\bm{J}$ & Matriz Jacobiana del mapeo isoparamétrico & m \\
$\det\bm{J}$ & Determinante del Jacobiano & m$^2$ \\
$\bm{B}$ & Matriz de deformación-desplazamiento & m$^{-1}$ \\
$\bm{D}$ & Matriz constitutiva (elástica) & Pa \\
$\bm{k}_e$ & Matriz de rigidez elemental & N/m \\
$\bm{K}$ & Matriz de rigidez global & N/m \\
$\bm{F}$ & Vector de cargas global & N \\
$\bm{R}$ & Vector de reacciones en los apoyos & N \\
$\bm{E}$ & Matriz de extrapolación de tensiones a los nodos & adimensional \\
$w_p$ & Peso de la cuadratura de Gauss & adimensional \\
$n_g$ & Número de puntos de Gauss del elemento & adimensional \\
$q_{SJ}$ & Jacobiano escalado (métrica de calidad, $\in[-1,1]$) & adimensional \\
$Q$ & Compacidad o \emph{stretch} (métrica de calidad, $\in[0,1]$) & adimensional \\
$L_{\min}$ & Longitud de la arista más corta del elemento & m \\
$D_{\max}$ & Longitud de la diagonal más larga del elemento & m \\
$\|\cdot\|_{L^2}$ & Norma $L^2$ del error en desplazamiento & según magnitud \\
$|\cdot|_{H^1}$ & Seminorma $H^1$ del error (gradiente) & según magnitud \\
\end{longtable}

\section*{Abreviaturas y siglas}

\begin{longtable}{@{}p{2.4cm}p{11.4cm}@{}}
\toprule
\textbf{Sigla} & \textbf{Significado} \\
\midrule
\endfirsthead
\toprule
\textbf{Sigla} & \textbf{Significado} \\
\midrule
\endhead
\bottomrule
\endfoot
MEF & Método de los Elementos Finitos \\
GDL & Grado(s) de libertad \\
TP & Tensión plana \\
DP & Deformación plana \\
Q4 & Elemento cuadrilátero isoparamétrico de 4 nodos (bilineal) \\
Q9 & Elemento cuadrilátero isoparamétrico de 9 nodos (bicuadrático) \\
MMS & Método de las Soluciones Manufacturadas \\
V\&V & Verificación y Validación \\
LU & Factorización LU (descomposición triangular inferior-superior) \\
COO & Formato disperso de coordenadas (\emph{coordinate list}) \\
CSR & Formato disperso de filas comprimidas (\emph{compressed sparse row}) \\
SRI & Integración reducida selectiva (\emph{selective reduced integration}) \\
MVC & Patrón de arquitectura Modelo-Vista-Controlador \\
GUI & Interfaz gráfica de usuario (\emph{graphical user interface}) \\
DXF & Formato de intercambio de dibujo CAD (\emph{drawing exchange format}) \\
CSV & Valores separados por comas (\emph{comma-separated values}) \\
PDF & Formato de documento portátil (\emph{portable document format}) \\
\end{longtable}
\renewcommand{\arraystretch}{1.0}

% Las dos longtable anteriores no llevan \caption, pero igualmente incrementan
% el contador `table`. Se reinicia para que la numeracion de tablas del cuerpo
% no quede desplazada (p. ej. que la primera tabla de la Introduccion no salga
% como "Tabla 3").
\setcounter{table}{0}
